% sec3_2.tex — Section 3.2: More Examples

Endogeneity bias arises in a wide variety of situations.  We examine a few more
examples.

%% ─── A Simple Macroeconomic Model ─────────────────────────────────────────
\subsection*{A Simple Macroeconomic Model}

Haavelmo's (1943) illustrative model is an extremely simple macroeconomic
model:
\begin{align*}
  C_i &= \alpha_0 + \alpha_1 Y_i + u_i, \quad 0 < \alpha_1 < 1
  \quad \text{(consumption function)} \\
  Y_i &= C_i + I_i \quad \text{(GNP identity)},
\end{align*}
where $C_i$ is aggregate consumption in year $i$, $Y_i$ is GNP, $I_i$ is investment, and
$\alpha_1$ is the Marginal Propensity to Consume out of income (the MPC).  As is familiar from
introductory macroeconomics, GNP in equilibrium is
\begin{equation}
  Y_i = \frac{\alpha_0}{1 - \alpha_1}
      + \frac{I_i}{1 - \alpha_1}
      + \frac{u_i}{1 - \alpha_1}.
  \label{eq:3.2.1}
\end{equation}
If investment is predetermined in that $\Cov(I_i, u_i) = 0$, it follows from \eqref{eq:3.2.1}
that
\[
  \Cov(Y_i, u_i) = \frac{\Var(u_i)}{1 - \alpha_1} > 0,
  \qquad
  \Cov(I_i, Y_i) = \frac{\Var(I_i)}{1 - \alpha_1} > 0.
\]
So income is endogenous in the consumption function, but investment is a valid
instrument for the endogenous regressor.  A straightforward calculation similar to
the one we used to derive \eqref{eq:3.1.7} shows that the OLS estimator of the MPC obtained
from regressing consumption on a constant and income is asymptotically biased:
\begin{equation}
  \plim\, \hat{\alpha}_{1,\text{OLS}} - \alpha_1
  = \frac{\Cov(Y_i, u_i)}{\Var(Y_i)}
  = \frac{1 - \alpha_1}{1 + \frac{\Var(I_i)}{\Var(u_i)}} > 0.
  \label{eq:3.2.2}
\end{equation}

This is perhaps the clearest example of simultaneity bias.  The difference from
Working's example of market equilibrium is that here the second equation is an
identity, which only makes it easier to see the endogeneity of the regressor.  The
asymptotic bias can be corrected for by the use of investment as the instrument
for income in the consumption function.  The role played by the observable supply
shifter in Working's example is played here by investment.

%% ─── Errors-in-Variables ──────────────────────────────────────────────────
\subsection*{Errors-in-Variables}

The term \textbf{errors-in-variables} refers to the phenomenon that an otherwise predetermined
regressor necessarily becomes endogenous when measured with error.  This
problem is ubiquitous, particularly in micro data on households.  For example, in
the Panel Study of Income Dynamics (PSID), information on variables such as
food consumption and income is collected over the telephone.  It is perhaps too
much to hope that the respondent can recall on the spot how much was spent on
food over a specified period of time.

The cross-section version of M.\ Friedman's (1957) Permanent Income Hypothesis
can be formulated neatly as an errors-in-variables problem.  The hypothesis
states that ``permanent consumption'' $C_i^*$ for household $i$ is proportional to ``permanent
income'' $Y_i^*$:
\begin{equation}
  C_i^* = k \, Y_i^* \quad \text{with } 0 < k < 1.
  \label{eq:3.2.3}
\end{equation}
It is assumed that measured consumption $C_i$ and measured income $Y_i$ are error-ridden
measures of permanent consumption and income:
\begin{equation}
  C_i = C_i^* + c_i \quad \text{and} \quad Y_i = Y_i^* + y_i.
  \label{eq:3.2.4}
\end{equation}
Measurement errors $c_i$ and $y_i$ are assumed to be zero mean and uncorrelated with
permanent consumption and income:
\begin{align}
  &\E(c_i) = 0, \quad \E(y_i) = 0, \quad \E(c_i y_i) = 0,
  \label{eq:3.2.5} \\
  &\E(C_i^* c_i) = 0, \quad \E(Y_i^* y_i) = 0, \quad
  \E(C_i^* y_i) = 0, \quad \E(Y_i^* c_i) = 0.
  \label{eq:3.2.6}
\end{align}
Substituting \eqref{eq:3.2.4} into \eqref{eq:3.2.3}, the relationship can be expressed in
terms of measured consumption and income:
\begin{equation}
  C_i = k \, Y_i + u_i \quad \text{with } u_i \equiv c_i - k \, y_i.
  \label{eq:3.2.7}
\end{equation}

This example differs from the previous ones in that the equation does not
include a constant.  So we should examine the cross moment $\E(Y_i u_i)$ rather than
the covariance $\Cov(Y_i, u_i)$ to determine whether income is predetermined.  It is
straightforward to derive from \eqref{eq:3.2.4}--\eqref{eq:3.2.7} that
\[
  \E(Y_i u_i) = -k\, \E(y_i^2) < 0.
\]
So measured income is endogenous in the consumption function \eqref{eq:3.2.7}.  Unlike
in the previous examples, the endogeneity is due to measurement errors.  Using
the fact that the OLS estimator of the $Y_i$ coefficient in \eqref{eq:3.2.7} is consistent for
the corresponding least squares projection coefficient $\E(C_i Y_i)/\E(Y_i^2)$, we can also
derive from \eqref{eq:3.2.4}--\eqref{eq:3.2.7} that
\begin{equation}
  \plim\, \hat{k}_{\text{OLS}}
  = \frac{k\, \E[(Y_i^*)^2]}{\E[(Y_i^*)^2] + \E(y_i^2)} < k.
  \label{eq:3.2.8}
\end{equation}
So the regression of measured consumption on measured income (without the intercept)
underestimates $k$.  Friedman used this result to explain why the MPC from
the cross-section regression of consumption on income is lower than the MPC estimated
from aggregate time-series data.

Let us suppose for a moment that there exists a valid instrument $x_i$.  So
$\E(x_i u_i) = 0$ and $\E(x_i Y_i) \neq 0$.  A similar argument we used for deriving
\eqref{eq:3.1.11} establishes that
\begin{equation}
  k = \frac{\E(x_i C_i)}{\E(x_i Y_i)}.
  \label{eq:3.2.9}
\end{equation}
So a consistent IV estimator is
\begin{equation}
  \hat{k}_{\text{IV}}
  = \frac{\text{sample mean of } x_i C_i}
         {\text{sample mean of } x_i Y_i}.
  \label{eq:3.2.10}
\end{equation}
But is there a valid instrument?  Yes, and it is a constant.  Substituting $x_i = 1$
into \eqref{eq:3.2.10}, we obtain a consistent estimator of $k$ which is the ratio of the sample
mean of measured consumption to the sample mean of measured income.  This is
how Friedman estimated $k$.

%% ─── Production Function ──────────────────────────────────────────────────
\subsection*{Production Function}

In many contexts, the error term includes factors that are observable to the economic
agent under study but unobservable to the econometrician.  Endogeneity
arises when regressors are decisions made by the agent on the basis of such factors.
Consider a cross-sectional sample of firms choosing labor input to maximize
profits.  The production function of the $i$-th firm is
\begin{equation}
  Q_i = A_i \cdot (L_i)^{\phi_1} \cdot \exp(v_i), \quad
  0 < \phi_1 < 1,
  \label{eq:3.2.11}
\end{equation}
where $Q_i$ is the firm's output, $L_i$ is labor input, $A_i$ is the firm's efficiency level
known to the firm, and $v_i$ is the technology shock.  In contrast to $A_i$, $v_i$ is not
observable to the firm when it chooses $L_i$.  Neither $A_i$ nor $v_i$ is observable to the
econometrician.

Assume that, for each firm $i$, $v_i$ is serially independent over time.  Therefore,
$B \equiv \E[\exp(v_i)]$ is the same for all $i$,\footnote{If $v_i$ for firm $i$ were correlated
over time, the firm would use past values of $v_i$ to forecast the current value of $v_i$
when choosing labor input, and so $B$ would differ across firms.  Also, since the expected
value of a nonlinear function of a random variable whose mean is zero is not necessarily
zero, $B$ is not necessarily zero even if $\E(v_i) = 0$.} and the level of output the firm
expects when it chooses $L_i$ is
\[
  A_i \cdot (L_i)^{\phi_1} \cdot B.
\]
Let $p$ be the output price and $w$ the wage rate.  For simplicity, assume that all the
firms are from the same competitive industry so that $p$ and $w$ are constant across
firms.  Firm $i$'s objective is to choose $L_i$ to maximize expected profits
\[
  p \cdot A_i \cdot (L_i)^{\phi_1} \cdot B - w \cdot L_i.
\]
Take the derivative of this with respect to $L_i$, set it equal to zero, and solve it for
$L_i$, to obtain the profit-maximizing labor input level:
\begin{equation}
  L_i = \left(\frac{w}{p}\right)^{\frac{1}{\phi_1 - 1}}
        (A_i \, B \phi_1)^{\frac{1}{1 - \phi_1}}.
  \label{eq:3.2.12}
\end{equation}

Let $u_i$ be firm $i$'s deviation from the mean log efficiency: $u_i \equiv \log(A_i) -
\E[\log(A_i)]$, and let $\phi_0 \equiv \E[\log(A_i)]$ (so $\E(u_i) = 0$ and
$A_i = \exp(\phi_0 + u_i)$).  Then,
the production function \eqref{eq:3.2.11} and the labor demand function \eqref{eq:3.2.12}
in logs can be written as
\begin{align}
  \log(Q_i) &= \phi_0 + \phi_1 \log(L_i) + (v_i + u_i),
  \label{eq:3.2.13} \\
  \log(L_i) &= \beta_0 + \frac{1}{1 - \phi_1}\, u_i,
  \label{eq:3.2.14}
\end{align}
where
\[
  \beta_0 = \frac{1}{\phi_1 - 1}\,[\log(w/p) - \phi_0 - \log(\phi_1 B)].
\]
Because all firms face the same prices, $\beta_0$ is constant across firms.
Equation~\eqref{eq:3.2.14} shows that, in the log-linear production function
\eqref{eq:3.2.13}, $\log(L_i)$ is an endogenous regressor positively correlated with the error
term $(v_i + u_i)$ through $u_i$.  Thus, the OLS estimator of $\phi_1$ in the estimation of the
log-linear production function confounds the contribution of output of $u_i$ with labor's
contribution.  This example illustrates yet another source of endogeneity: a variable
chosen by the agent taking into account some error component unobservable to the
econometrician.

%% ─── Questions for Review ──────────────────────────────────────────────────
\bigskip
\begin{center}\rule{0.6\textwidth}{0.4pt}\end{center}
\noindent\textsc{Questions for Review}
\medskip

\begin{enumerate}
\item In Friedman's Permanent Income Hypothesis, consider the regression of $C_i$ on
  a constant and $Y_i$.  Derive the plim of the OLS estimator of the $Y_i$ coefficient in
  this regression with a constant.  \textbf{Hint:} The plim equals the ratio of $\Cov(C_i, Y_i)$
  to $\Var(Y_i)$.  Show that it equals
  \[
    \frac{k\,\Var(Y_i^*)}{\Var(Y_i^*) + \Var(y_i)}.
  \]

\item In the production function example, show that
  $\plim_{n\to\infty} \hat{\phi}_{1,\text{OLS}} = 1$, where
  $\hat{\phi}_{1,\text{OLS}}$ is the OLS estimate of $\phi_1$ from \eqref{eq:3.2.13}.
  \textbf{Hint:} Eliminate $u_i$ from the log output equation \eqref{eq:3.2.13} by using the
  labor demand equation \eqref{eq:3.2.14}.

\item In the production function example, suppose the firm can observe $v_i$ as well
  as $u_i$ before deciding on labor input.  How does the demand equation for labor
\end{enumerate}
